% TEMPLATE-EXAMEN-MILC-v3.tex - plantilla maestra para los exámenes finales
% MILC v3 de Plataforma Conéctate (grados 8°-11°).
%
% Ver scripts/build-examenes.py para la lista completa de claves esperadas.
% El script valida que no queden placeholders sin reemplazar antes de
% escribir el archivo final.
%
% Diferencias clave respecto a template-milc-v3-conectate.tex (guías):
%   - Portada NEGRA (#000000), texto blanco. Coherente con el card del
%     bento del periodo en la web.
%   - Solo 5 preguntas (concreto · contexto · práctico ×2 · reflexivo).
%   - Cabecera con clave examen y "Pregunta N de 5".
%   - Espacio para respuesta del estudiante después de cada ítem.
%   - Aprendizajes cubiertos visibles al inicio.

\documentclass[11pt,letterpaper]{article}
\usepackage[margin=1.75cm]{geometry}
\usepackage[table]{xcolor}
\usepackage{fontspec}
\usepackage[spanish]{babel}
\usepackage{tikz}
\usepackage[most]{tcolorbox}
\usepackage{tabularx}
\usepackage{array}
\usepackage{fancyhdr}
\usepackage{titlesec}
\usepackage{ragged2e}
\usepackage{enumitem}
\usepackage{microtype}

\setmainfont{Helvetica Neue}
\setsansfont{Helvetica Neue}
\setlength{\parindent}{0pt}
\setlength{\parskip}{6pt}
\hyphenpenalty=200
\exhyphenpenalty=200
\pretolerance=400
\tolerance=2000
\setlength{\emergencystretch}{1em}
\renewcommand{\arraystretch}{1.25}
\setlist{leftmargin=5mm,itemsep=1mm,topsep=1mm}
\newcolumntype{Y}{>{\RaggedRight\arraybackslash}X}

% Paleta institucional MILC (compartida con guías).
\definecolor{milcMagenta}{HTML}{E6007E}
\definecolor{milcVino}{HTML}{5A0038}
\definecolor{milcTurquesa}{HTML}{0093A5}
\definecolor{milcVerde}{HTML}{58A923}
\definecolor{milcMostaza}{HTML}{E5B400}
\definecolor{milcOcre}{HTML}{B86600}
\definecolor{milcNegro}{HTML}{241816}
\definecolor{milcGris}{HTML}{F7F7F4}
\definecolor{milcLinea}{HTML}{E8E2DD}

% Paleta del PERIODO (chip de color en portada + acentos).
\definecolor{milcPeriodoPrimary}{HTML}{<<<COLOR_PERIODO_PRIMARY>>>}
\definecolor{milcPeriodoDark}{HTML}{<<<COLOR_PERIODO_DARK>>>}
\definecolor{milcPeriodoSoft}{HTML}{<<<COLOR_PERIODO_SOFT>>>}

% Fondo negro absoluto para la portada (coherente con bento card web).
\definecolor{examenFondo}{HTML}{000000}
\definecolor{examenTexto}{HTML}{FFFFFF}
\definecolor{examenTextoSubtle}{HTML}{B8B8B8}
\color{milcNegro}

\pagestyle{fancy}
\fancyhf{}
\lhead{\small Examen final · Tecnología e Informática · Grado <<<GRADO>>>}
\rhead{\small Examen <<<CLAVE>>> · MILC v3}
\cfoot{\small\thepage}
\renewcommand{\headrulewidth}{0pt}

% ─── Macros visuales ────────────────────────────────────────────────

\newtcolorbox{softbox}[3]{
  enhanced,colback=#2,colframe=#1,borderline west={4pt}{0pt}{#1},
  arc=2mm,boxrule=.4pt,left=4mm,right=4mm,top=3mm,bottom=3mm,
  before skip=6pt,after skip=8pt,title={#3},coltitle=white,
  colbacktitle=#1,fonttitle=\sffamily\bfseries,fontupper=\RaggedRight,
  attach boxed title to top left={xshift=3mm,yshift=-2mm},
  boxed title style={arc=2mm, boxrule=0pt}
}

\newtcolorbox{infoband}[1]{
  enhanced,colback=#1!8,colframe=#1!50,arc=2mm,boxrule=.5pt,
  left=4mm,right=4mm,top=2mm,bottom=2mm,before skip=4pt,after skip=4pt,
  fontupper=\small\RaggedRight
}

% Cabecera de pregunta: número en chip + título imperativo.
\newtcolorbox{cabecerapregunta}[2]{
  enhanced,colback=milcPeriodoSoft,colframe=milcPeriodoDark,
  borderline west={4pt}{0pt}{milcPeriodoDark},
  arc=2mm,boxrule=0pt,left=5mm,right=4mm,top=3mm,bottom=3mm,
  before skip=10pt,after skip=4pt,
  title={#1},coltitle=white,
  colbacktitle=milcPeriodoDark,fonttitle=\sffamily\bfseries\large,
  fontupper=\sffamily\bfseries\large\color{milcNegro},
  attach boxed title to top left={xshift=4mm,yshift=-3mm},
  boxed title style={arc=2mm, boxrule=0pt}
}

% Caja "voces" para el ítem reflexivo (las 3 lentes visibles).
\newtcolorbox{vozbox}[2]{
  enhanced,colback=white,colframe=#1,
  arc=2mm,boxrule=.5pt,left=4mm,right=4mm,top=2mm,bottom=2mm,
  before skip=4pt,after skip=4pt,
  title={#2},coltitle=white,colbacktitle=#1,
  fonttitle=\sffamily\bfseries\small,
  fontupper=\itshape\small\RaggedRight,
  attach boxed title to top left={xshift=3mm,yshift=-2mm},
  boxed title style={arc=2mm, boxrule=0pt}
}

% Líneas para que el estudiante escriba sus respuestas.
\newcommand{\respuesta}[1]{%
  \par\vspace{2mm}%
  \foreach \n in {1,...,#1}{%
    \noindent\color{milcLinea}\rule{\linewidth}{0.5pt}\par\vspace{5mm}%
  }%
  \color{milcNegro}%
}

% ─── Portada ─────────────────────────────────────────────────────────

\begin{document}
\thispagestyle{empty}
\begin{tikzpicture}[remember picture,overlay]
  % Fondo negro absoluto (toda la página).
  \fill[examenFondo] (current page.south west) rectangle (current page.north east);

  % Chip institucional arriba a la izquierda.
  \node[anchor=north west,text=examenTexto] at ([xshift=2.0cm,yshift=-1.85cm]current page.north west)
    {\sffamily\bfseries\fontsize{11}{13}\selectfont EXAMEN FINAL};
  \node[anchor=north west,text=examenTextoSubtle] at ([xshift=2.0cm,yshift=-2.45cm]current page.north west)
    {\sffamily\bfseries\fontsize{9}{11}\selectfont GRADO <<<GRADO>>>° · PERÍODO <<<PERIODO>>> · TECNOLOGÍA E INFORMÁTICA};

  % Chip de color del periodo arriba a la derecha (acento sutil).
  \begin{scope}[shift={([xshift=-2.5cm,yshift=-2.15cm]current page.north east)}]
    \fill[milcPeriodoPrimary,rounded corners=2mm] (-1.5,-.4) rectangle (1.5,.4);
    \node[text=white] at (0,0) {\sffamily\bfseries\fontsize{9}{11}\selectfont APLICACIÓN · SESIÓN 10};
  \end{scope}

  % Título principal grande (centrado, blanco).
  \node[anchor=west,text=examenTexto,text width=17cm,align=left] at ([xshift=2.0cm,yshift=-9.5cm]current page.north west)
    {\sffamily\bfseries\fontsize{36}{42}\selectfont <<<TITULO_PORTADA>>>};

  % Subtítulo (gris claro).
  \node[anchor=west,text=examenTextoSubtle,text width=17cm,align=left] at ([xshift=2.0cm,yshift=-14.5cm]current page.north west)
    {\sffamily\fontsize{14}{18}\selectfont Síntesis · Aplicación · Argumentación.\\Verifica la transferencia de lo aprendido en las 10 sesiones del período.};

  % Pie con metadatos (duración, fecha, ponderación, autor).
  \path[draw=examenTextoSubtle,line width=.4pt]
    ([xshift=2.0cm,yshift=4.0cm]current page.south west) --
    ([xshift=-2.0cm,yshift=4.0cm]current page.south east);

  \node[anchor=west,text=examenTexto,text width=17cm,align=left] at ([xshift=2.0cm,yshift=3.0cm]current page.south west)
    {\sffamily\bfseries\fontsize{11}{14}\selectfont DURACIÓN \quad\textcolor{examenTextoSubtle}{<<<DURACION_MIN>>> min} \hspace{1.0cm} APLICACIÓN \quad\textcolor{examenTextoSubtle}{<<<FECHA_APLICACION>>>} \hspace{1.0cm} PONDERACIÓN \quad\textcolor{examenTextoSubtle}{<<<PONDERACION>>>\%}};

  \node[anchor=west,text=examenTextoSubtle,text width=17cm,align=left] at ([xshift=2.0cm,yshift=1.9cm]current page.south west)
    {\sffamily\fontsize{10}{12}\selectfont Institución Educativa Sor María Juliana · Cartago, Valle del Cauca\\Autor: Dr. Álvaro Cárdenas Orozco · Metodología MILC v3};
\end{tikzpicture}
\mbox{}
\newpage

% ─── Datos del estudiante (página 2) ─────────────────────────────────

\section*{Datos del estudiante}

\noindent
\begin{tabularx}{\linewidth}{@{}p{3.5cm}Y@{}}
Nombre completo: & \color{milcLinea}\rule{\linewidth}{0.5pt}\color{milcNegro}\\[3mm]
Grado y grupo: & \color{milcLinea}\rule{\linewidth}{0.5pt}\color{milcNegro}\\[3mm]
Fecha: & \color{milcLinea}\rule{\linewidth}{0.5pt}\color{milcNegro}\\[3mm]
Firma: & \color{milcLinea}\rule{\linewidth}{0.5pt}\color{milcNegro}\\
\end{tabularx}

\vspace{4mm}

% ─── Instrucciones ───────────────────────────────────────────────────

\begin{softbox}{milcPeriodoDark}{milcPeriodoSoft}{Instrucciones}
<<<INSTRUCCIONES>>>
\end{softbox}

% ─── Anclaje MILC: contexto + saber ancestral ────────────────────────

\begin{softbox}{milcVino}{milcGris}{Saber ancestral del período}
<<<SABER_ANCESTRAL>>>
\end{softbox}

\begin{softbox}{milcTurquesa}{milcGris}{Contexto del examen}
<<<CONTEXTO>>>
\end{softbox}

\begin{softbox}{milcMostaza}{milcGris}{Pregunta-marco}
\textit{<<<PREGUNTA_MARCO>>>}
\end{softbox}

% ─── Aprendizajes cubiertos ──────────────────────────────────────────

\section*{Aprendizajes que entran al examen}

\begin{infoband}{milcPeriodoDark}
<<<APRENDIZAJES_LISTA>>>
\end{infoband}

\newpage

% ─── Pregunta 1 · Concreto (¿Qué aprendí?) ───────────────────────────

\begin{cabecerapregunta}{Pregunta 1 de 5 · Concreto}{¿Qué aprendí?}
<<<CONCRETO_ENUNCIADO>>>
\end{cabecerapregunta}

\begin{infoband}{milcVerde}
\textbf{Criterio.} <<<CONCRETO_CRITERIO>>>
\end{infoband}

\respuesta{6}

\newpage

% ─── Pregunta 2 · Contexto (saber ancestral) ─────────────────────────

\begin{cabecerapregunta}{Pregunta 2 de 5 · Contexto}{Conecta el oficio con el código.}
<<<CONTEXTO_ENUNCIADO>>>
\end{cabecerapregunta}

\begin{infoband}{milcVerde}
\textbf{Criterio.} <<<CONTEXTO_CRITERIO>>>
\end{infoband}

\respuesta{6}

\newpage

% ─── Pregunta 3 · Práctico 1 ─────────────────────────────────────────

\begin{cabecerapregunta}{Pregunta 3 de 5 · Práctico}{Aplica lo que aprendiste.}
<<<PRACTICO_1_ENUNCIADO>>>
\end{cabecerapregunta}

<<<PRACTICO_1_CONTEXTO_BLOQUE>>>

<<<PRACTICO_1_OPCIONES_BLOQUE>>>

\begin{infoband}{milcVerde}
\textbf{Criterio.} <<<PRACTICO_1_CRITERIO>>>
\end{infoband}

\respuesta{8}

\newpage

% ─── Pregunta 4 · Práctico 2 ─────────────────────────────────────────

\begin{cabecerapregunta}{Pregunta 4 de 5 · Práctico}{Aplica lo que aprendiste.}
<<<PRACTICO_2_ENUNCIADO>>>
\end{cabecerapregunta}

<<<PRACTICO_2_CONTEXTO_BLOQUE>>>

<<<PRACTICO_2_OPCIONES_BLOQUE>>>

\begin{infoband}{milcVerde}
\textbf{Criterio.} <<<PRACTICO_2_CRITERIO>>>
\end{infoband}

\respuesta{8}

\newpage

% ─── Pregunta 5 · Reflexivo (triángulo de pensamiento) ───────────────

\begin{cabecerapregunta}{Pregunta 5 de 5 · Reflexivo}{Aplica el triángulo de pensamiento.}
<<<REFLEXIVO_ENUNCIADO>>>
\end{cabecerapregunta}

\begin{vozbox}{milcVino}{Dussel · lente del nosotros}
<<<VOZ_DUSSEL>>>
\end{vozbox}

\begin{vozbox}{milcMostaza}{Estoico · lente del cuidado interior}
<<<VOZ_ESTOICO>>>
\end{vozbox}

\begin{vozbox}{milcTurquesa}{Floridi · lente de la infoesfera}
<<<VOZ_FLORIDI>>>
\end{vozbox}

\begin{infoband}{milcVerde}
\textbf{Criterio.} <<<REFLEXIVO_CRITERIO>>>
\end{infoband}

\respuesta{8}

\newpage

% ─── Cierre ──────────────────────────────────────────────────────────

\section*{Cierre del período}

\begin{softbox}{milcPeriodoDark}{milcPeriodoSoft}{Una última invitación}
<<<CIERRE>>>
\end{softbox}

\vspace{8mm}

\begin{center}
\textcolor{milcPeriodoDark}{\sffamily\bfseries\large
Institución Educativa Sor María Juliana
}\\[1mm]
\textcolor{milcNegro}{\sffamily\small
Cartago, Valle del Cauca · Tecnología e Informática · Dr. Álvaro Cárdenas Orozco
}
\end{center}

\end{document}
